% appendix.tex -- Supplementary Material

\section{Sequence Encoding, PCA Projection, and Decoding}\label{app:pca}

Stochastic attention (SA) operates on continuous vectors on the unit sphere,
while protein sequences are discrete objects over a 20-letter amino acid
alphabet. Bridging the two requires three transformations---one-hot encoding,
PCA dimensionality reduction, and unit-norm projection---and the reverse path
(inverse PCA followed by argmax decoding) recovers discrete sequences from
continuous samples. We describe each step below, define all notation, and then
discuss how the multiplicity vector $\mathbf{r}$ (the per-pattern weighting
used for functional conditioning; see Section~\ref{sec:theory}) interacts with
the PCA representation.

\paragraph{One-hot encoding.}
Given a cleaned alignment of $K$ sequences, each of length $L$ amino acids
over the standard 20-letter alphabet
$\mathcal{A} = \{\text{A}, \text{R}, \text{N}, \ldots, \text{V}\}$, we
encoded each sequence as a binary vector in $\mathbb{R}^{20L}$ by
concatenating per-position indicator vectors. For sequence $k$ ($k = 1,
\ldots, K$) with amino acid $a_\ell^{(k)}$ at alignment position $\ell$
($\ell = 1, \ldots, L$):
\begin{equation}
  \mathbf{x}_k = \bigl[\mathbf{e}_{a_1^{(k)}},\; \mathbf{e}_{a_2^{(k)}},\;
    \ldots,\; \mathbf{e}_{a_L^{(k)}}\bigr]^\top \in \{0,1\}^{20L},
\end{equation}
where $\mathbf{e}_a \in \{0,1\}^{20}$ is the standard basis vector for amino
acid $a$ (i.e., a one-hot indicator with a 1 in the position corresponding to
$a$ and 0 elsewhere). The full one-hot dimensionality is
$d_{\mathrm{full}} = 20L$, which is the total number of coordinates in the
binary representation. For the four families studied here,
$d_{\mathrm{full}}$ ranged from 520 ($\omega$-conotoxin, $L{=}26$) to 1{,}060
(Kunitz, $L{=}53$).

\paragraph{Why dimensionality reduction is necessary.}
The full one-hot space is vastly over-dimensioned relative to the number of
stored patterns $K$. The ratio $d_{\mathrm{full}}/K$ ranged from 7.0
($\omega$-conotoxin, $K{=}74$, $d_{\mathrm{full}}{=}520$) to 10.7 (Kunitz,
$K{=}99$, $d_{\mathrm{full}}{=}1{,}060$) across the four families. This
over-dimensioning creates a problem for the modern Hopfield
energy~\cite{ramsauerHopfieldNetworksAll2021}: the similarity score between a
stored pattern $\mathbf{m}_k$ (a unit-norm vector representing sequence $k$;
defined below) and a query state $\boldsymbol{\xi}$ (the current position of
the Langevin sampler) is the dot product
$e_k = \mathbf{m}_k^\top\boldsymbol{\xi}$. For a random query drawn uniformly
from the unit sphere $\mathbb{S}^{d-1}$ (the surface of the $d$-dimensional
unit ball), these scores have variance $\mathrm{Var}(e_k) = 1/d$. When $d$ is
large, all scores concentrate near zero and the inverse temperature $\beta$
(a scalar parameter controlling the sharpness of the softmax attention) must be
set extremely high to differentiate among patterns, making Langevin sampling
inefficient. PCA removes this redundancy by projecting onto the $d \ll
d_{\mathrm{full}}$ directions of actual variation in the alignment.

\paragraph{PCA projection.}
We computed PCA from the one-hot encoded alignment. Let
$\mathbf{X} = [\mathbf{x}_1, \ldots, \mathbf{x}_K] \in
\mathbb{R}^{d_{\mathrm{full}} \times K}$ be the matrix whose columns are the
one-hot encoded sequences. We first centered the data by subtracting the
column mean:
\begin{equation}
  \bar{\mathbf{x}} = \frac{1}{K}\sum_{k=1}^K \mathbf{x}_k \in
  \mathbb{R}^{d_{\mathrm{full}}}, \qquad
  \tilde{\mathbf{X}} = \mathbf{X} - \bar{\mathbf{x}}\mathbf{1}_K^\top,
\end{equation}
where $\bar{\mathbf{x}}$ is the mean one-hot vector (i.e., the family's
position-specific amino acid frequency profile) and $\mathbf{1}_K$ is the
$K$-vector of ones. We then computed the economy singular value decomposition
(SVD):
\begin{equation}
  \tilde{\mathbf{X}} = \mathbf{U}\boldsymbol{\Sigma}\mathbf{V}^\top,
\end{equation}
where $\mathbf{U} \in \mathbb{R}^{d_{\mathrm{full}} \times r}$ contains the
left singular vectors (the principal directions of variation in
one-hot space), $\boldsymbol{\Sigma} = \mathrm{diag}(\sigma_1, \ldots,
\sigma_r)$ contains the singular values in decreasing order
($\sigma_1 \geq \cdots \geq \sigma_r > 0$),
$\mathbf{V} \in \mathbb{R}^{K \times r}$ contains the right singular vectors,
and $r = \mathrm{rank}(\tilde{\mathbf{X}}) \leq \min(d_{\mathrm{full}}, K) - 1$
is the rank of the centered data matrix. The fraction of total variance
explained by the first $p$ principal components is:
\begin{equation}
  \rho_{\mathrm{var}}(p) = \frac{\sum_{j=1}^p \sigma_j^2}{\sum_{j=1}^r \sigma_j^2}.
\end{equation}
We selected the retained PCA dimension $d$ as the smallest integer $p$ such
that $\rho_{\mathrm{var}}(p) \geq 0.95$, i.e., retaining at least 95\% of
the total variance. Across the four families, $d$ ranged from 20
($\omega$-conotoxin) to 186 (WW), yielding compression ratios
$d_{\mathrm{full}}/d$ of $3.3\times$ (WW) to $26\times$
($\omega$-conotoxin) (Table~\ref{tab:cross-family}). The projection matrix is
$\mathbf{W}_d = \mathbf{U}_{:,1:d} \in
\mathbb{R}^{d_{\mathrm{full}} \times d}$ (the first $d$ columns of
$\mathbf{U}$), and each sequence is projected as:
\begin{equation}
  \mathbf{z}_k = \mathbf{W}_d^\top(\mathbf{x}_k - \bar{\mathbf{x}}) \in \mathbb{R}^d,
\end{equation}
where $\mathbf{z}_k$ is the PCA-reduced representation of sequence $k$.

\paragraph{Unit-norm projection.}
After PCA, we normalized each projected vector to unit $L_2$ norm:
\begin{equation}
  \mathbf{m}_k = \frac{\mathbf{z}_k}{\|\mathbf{z}_k\|_2} \in \mathbb{S}^{d-1},
\end{equation}
where $\mathbb{S}^{d-1} = \{\mathbf{v} \in \mathbb{R}^d : \|\mathbf{v}\|_2 = 1\}$
is the $(d{-}1)$-dimensional unit sphere. This normalization ensures that
similarity scores $e_k = \mathbf{m}_k^\top\boldsymbol{\xi}$ lie in $[-1, 1]$
and that the concentration-of-measure result $\mathrm{Var}(e_k) = 1/d$ holds
for random queries. It also prevents patterns with larger PCA norms (sequences
that deviate more from the family mean) from dominating the softmax attention
weights by virtue of their norm rather than their direction. The columns
$\mathbf{m}_1, \ldots, \mathbf{m}_K$ form the memory matrix
$\hat{\mathbf{X}} = [\mathbf{m}_1, \ldots, \mathbf{m}_K]
\in \mathbb{R}^{d \times K}$ used throughout the paper.

\paragraph{Decoding: from continuous samples to amino acid sequences.}
The Langevin sampler (Eq.~\ref{eq:ula}) produces continuous state vectors
$\boldsymbol{\xi} \in \mathbb{R}^d$ in the $d$-dimensional PCA space. To
decode a sample back to an amino acid sequence, we applied two steps. First,
inverse PCA reconstruction maps from PCA space back to the full one-hot space:
\begin{equation}
  \hat{\mathbf{x}} = \bar{\mathbf{x}} + \mathbf{W}_d\boldsymbol{\xi}
  \in \mathbb{R}^{d_{\mathrm{full}}},
\end{equation}
where $\hat{\mathbf{x}}$ is a continuous vector that approximates a one-hot
encoding but is not exactly binary. Second, per-position argmax decoding
assigns the amino acid with the largest coordinate at each alignment position:
\begin{equation}
  a_\ell = \arg\max_{a \in \mathcal{A}}\; \hat{x}_{20(\ell-1) + a},
  \qquad \ell = 1, \ldots, L.
\end{equation}
This decoding is deterministic and produces a valid amino acid at every
position. In practice, 100\% of decoded positions yielded standard amino acids
across all four families and all conditions tested.

\paragraph{Effect of the multiplicity vector $\mathbf{r}$ on the PCA
representation.}
A key property of multiplicity-weighted conditioning is that the multiplicity
vector $\mathbf{r} = (r_1, \ldots, r_K)^\top$ (where $r_k > 0$ is the weight
assigned to stored pattern $k$; see Section~\ref{sec:theory}) does \emph{not}
alter the PCA encoding. The memory matrix $\hat{\mathbf{X}}$ is constructed
once from the full family alignment and remains fixed regardless of the
multiplicity ratio $\rho = r_{\mathrm{des}}/r_{\mathrm{bg}}$ (the ratio of
weights assigned to designated versus background patterns). The multiplicity
weights enter only through the softmax logits: the weighted score function is
\begin{equation}
  \nabla_{\boldsymbol{\xi}}\log p_{\mathbf{r}}(\boldsymbol{\xi})
  = \hat{\mathbf{X}}\,\mathrm{softmax}(\beta\,\hat{\mathbf{X}}^\top
  \boldsymbol{\xi} + \log\mathbf{r}) - \boldsymbol{\xi},
\end{equation}
where $\log\mathbf{r} = (\log r_1, \ldots, \log r_K)^\top$ is the
elementwise log-multiplicity vector that is added to the pre-softmax logits
(Proposition~\ref{prop:score}). This separation of PCA encoding from
multiplicity weighting has three consequences:

\begin{enumerate}
  \item \textbf{The PCA subspace is independent of $\rho$.}
  The principal directions $\mathbf{W}_d$, the centering vector
  $\bar{\mathbf{x}}$, and the retained dimension $d$ are all determined by the
  alignment alone. Changing $\rho$ does not add, remove, or reweight any
  principal component. This means the decoding pipeline---and hence the
  calibration gap $\Delta = f_{\mathrm{eff}} - f_{\mathrm{obs}}$ (the
  difference between the intended and observed phenotype fraction; see
  Section~\ref{sec:gap})---is identical for all $\rho$ values on a given
  family.

  \item \textbf{$\rho$ controls attention, not geometry.}
  Increasing $\rho$ shifts attention weight toward designated patterns by
  adding $\log\rho$ to their pre-softmax logits. We confirmed empirically that
  the mean attention weight on designated patterns
  $\bar{a}_{\mathrm{des}}$ (the average fraction of the softmax probability
  mass assigned to designated patterns during sampling) tracked the effective
  designated fraction $f_{\mathrm{eff}}(\rho) = K_{\mathrm{des}}\rho /
  (K_{\mathrm{des}}\rho + K_{\mathrm{bg}})$ with $< 0.3\%$ deviation across
  all $\rho$ values and all families (Table~\ref{fig:per-family-rho}).
  However, this attention shift must still pass through the PCA reconstruction
  and argmax decoding to produce discrete residues, which is where the
  calibration gap arises.

  \item \textbf{$\rho$ shifts the phase transition.}
  The logit biases $\log r_k$ pre-concentrate the attention distribution even
  at $\beta = 0$: the weighted attention entropy at zero inverse temperature is
  $H_{\mathbf{r}}(0) = \log K_{\mathrm{eff}}(\mathbf{r}) < \log K$, where
  $K_{\mathrm{eff}}(\mathbf{r}) = (\sum_k r_k)^2 / \sum_k r_k^2$ is the
  effective number of patterns (a measure of how many patterns carry
  significant weight under $\mathbf{r}$; it equals $K$ when all weights are
  equal and converges to $K_{\mathrm{des}}$ as $\rho \to \infty$). Because the
  starting entropy is lower, reaching the retrieval-phase inflection requires a
  higher inverse temperature, so $\beta^{*}(\rho) \geq \beta^{*}(1)$, where
  $\beta^{*}$ is the critical inverse temperature at which SA transitions from
  an exploration phase (attention spread across many patterns) to a retrieval
  phase (attention concentrated on one or a few patterns). We observed this
  empirically: $\beta^{*}$ increased from 4.3 ($\rho=1$,
  $K_{\mathrm{eff}} = 99$) to 10.1 ($\rho=1{,}000$,
  $K_{\mathrm{eff}} = 32.1$) on the Kunitz domain
  (Fig.~\ref{fig:phase-transition}).
\end{enumerate}

By contrast, hard curation---where the memory matrix is built from the
designated subset only---\emph{does} change the PCA. Under hard curation, the
centering vector $\bar{\mathbf{x}}$, the principal directions
$\mathbf{W}_d$, and the retained dimension $d$ are all computed from the
$K_{\mathrm{des}}$ designated sequences rather than the full $K$-sequence
alignment. The curated PCA therefore captures subset-specific variation that
may be orthogonal to the leading components of the full-family PCA. This is
why hard curation achieved 100\% phenotype transfer regardless of the Fisher
separation index $S$ (a measure of how well the full-family PCA separates
designated from background sequences; see Section~\ref{sec:gap}), while
multiplicity weighting was limited by the full-family PCA geometry.

\section{Supplementary Tables}

\subsection*{Table S1: Full per-family $\rho$ sweep data}

Selected values from the $\rho$ sweep on each of the three benchmark families.
$f_{\mathrm{eff}}$ is the effective designated fraction computed from the multiplicity
vector; $f_{\mathrm{obs}}$ is the hard-decoded P1/marker fraction averaged over 930
generated sequences; $\bar{a}_{\mathrm{des}}$ is the mean softmax attention weight on
designated patterns; $D$ is mean pairwise diversity.

\begin{table}[h]
  \centering
  \caption{Kunitz domain (PF00014) $\rho$ sweep. $K=99$, $K_{\mathrm{des}}=32$.}
  \label{tab:kunitz-rho}
  \small
  \begin{tabular}{ccccc}
    \toprule
    $\rho$ & $f_{\mathrm{eff}}$ & $f_{\mathrm{obs}}$ & $\bar{a}_{\mathrm{des}}$ & $D$ \\
    \midrule
    1    & 0.323 & 0.406 & 0.324 & 0.563 \\
    5    & 0.705 & 0.481 & 0.712 & 0.556 \\
    10   & 0.827 & 0.539 & 0.832 & 0.547 \\
    50   & 0.962 & 0.598 & 0.963 & 0.526 \\
    100  & 0.980 & 0.615 & 0.981 & 0.520 \\
    200  & 0.990 & 0.623 & 0.991 & 0.519 \\
    500  & 0.996 & 0.631 & 0.997 & 0.515 \\
    1000 & 0.998 & 0.634 & 0.998 & 0.514 \\
    \bottomrule
  \end{tabular}
\end{table}

\begin{table}[h]
  \centering
  \caption{SH3 domain (PF00018) $\rho$ sweep. $K=55$, $K_{\mathrm{des}}=33$.}
  \label{tab:sh3-rho}
  \small
  \begin{tabular}{ccccc}
    \toprule
    $\rho$ & $f_{\mathrm{eff}}$ & $f_{\mathrm{obs}}$ & $\bar{a}_{\mathrm{des}}$ & $D$ \\
    \midrule
    1    & 0.600 & 0.852 & 0.614 & 0.580 \\
    5    & 0.882 & 0.897 & 0.890 & 0.575 \\
    10   & 0.938 & 0.906 & 0.941 & 0.575 \\
    50   & 0.983 & 0.965 & 0.983 & 0.567 \\
    100  & 0.991 & 0.977 & 0.991 & 0.559 \\
    200  & 0.996 & 0.984 & 0.996 & 0.554 \\
    500  & 0.999 & 0.989 & 0.999 & 0.513 \\
    1000 & 0.999 & 0.992 & 0.999 & 0.510 \\
    \bottomrule
  \end{tabular}
\end{table}

\begin{table}[h]
  \centering
  \caption{WW domain (PF00397) $\rho$ sweep. $K=420$, $K_{\mathrm{des}}=69$.}
  \label{tab:ww-rho}
  \small
  \begin{tabular}{ccccc}
    \toprule
    $\rho$ & $f_{\mathrm{eff}}$ & $f_{\mathrm{obs}}$ & $\bar{a}_{\mathrm{des}}$ & $D$ \\
    \midrule
    1    & 0.164 & 0.131 & 0.152 & 0.690 \\
    5    & 0.496 & 0.184 & 0.487 & 0.689 \\
    10   & 0.663 & 0.226 & 0.661 & 0.676 \\
    50   & 0.901 & 0.277 & 0.901 & 0.657 \\
    100  & 0.947 & 0.303 & 0.947 & 0.649 \\
    200  & 0.973 & 0.327 & 0.974 & 0.643 \\
    500  & 0.990 & 0.360 & 0.991 & 0.642 \\
    1000 & 0.995 & 0.371 & 0.995 & 0.642 \\
    \bottomrule
  \end{tabular}
\end{table}

\subsection*{Table S2: $\beta$ sweep at fixed $\rho$ on Kunitz}

Phenotype fraction and diversity as a function of $\beta/\beta^{*}$ for three fixed $\rho$
values on the Kunitz domain. $f_{\mathrm{obs}}$ and $D$ are averages over 930 generated
sequences. All conditions produced 100\% valid residues.

\begin{table}[h]
  \centering
  \caption{$\beta$ sweep on Kunitz domain ($K=99$, $K_{\mathrm{des}}=32$).}
  \label{tab:beta-sweep}
  \small
  \begin{tabular}{cccc}
    \toprule
    $\rho$ & $\beta/\beta^{*}$ & $f_{\mathrm{obs}}$ & $D$ \\
    \midrule
    10  & 0.50 & 0.471 & 0.581 \\
    10  & 1.00 & 0.539 & 0.547 \\
    10  & 1.50 & 0.573 & 0.525 \\
    10  & 2.00 & 0.597 & 0.508 \\
    10  & 3.00 & 0.613 & 0.484 \\
    \midrule
    50  & 0.50 & 0.506 & 0.573 \\
    50  & 1.00 & 0.582 & 0.537 \\
    50  & 1.50 & 0.631 & 0.514 \\
    50  & 2.00 & 0.653 & 0.499 \\
    50  & 3.00 & 0.700 & 0.469 \\
    \midrule
    200 & 0.50 & 0.535 & 0.564 \\
    200 & 1.00 & 0.613 & 0.527 \\
    200 & 1.50 & 0.648 & 0.503 \\
    200 & 2.00 & 0.677 & 0.488 \\
    200 & 3.00 & 0.706 & 0.463 \\
    \bottomrule
  \end{tabular}
\end{table}

\section{Supplementary Figures}

Binding-loop per-position amino acid frequency heatmaps for the $\omega$-conotoxin experiment
are shown in Fig.~\ref{fig:conotoxin-loop} of the main text (positions 9--17 of the alignment).
Additional SI figures for the cross-family PCA separation analysis and full entropy-curve
families are described below.

\paragraph{PCA separation histograms.}
For each family, designated and background sequences were projected onto their top two principal
components (from the full-family PCA). The separation index $S$ was computed as the ratio of
the inter-class distance to the intra-class spread (see Methods). Kunitz
($S=0.20$) shows partial overlap; SH3 ($S=0.34$) shows clear separation; WW ($S=0.11$) shows
extensive overlap, consistent with the large calibration gap.

\paragraph{pLDDT versus TM-score correlation.}
Figure~\ref{fig:plddt-tmscore-scatter} shows the joint distribution of ESMFold pLDDT and
TM-score for all 250 predicted structures. SA-generated sequences cluster tightly in the
high-pLDDT, high-TM-score region (upper right), while HMM-emitted sequences scatter broadly
across lower values. The two metrics are positively correlated, confirming
that ESMFold confidence is a reliable proxy for structural accuracy on this family.

\paragraph{Entropy-curve families across $\rho$.}
Figure~\ref{fig:phase-transition} in the main text shows entropy curves for five representative
$\rho$ values on the Kunitz family. The full sweep ($\rho \in \{1, 2, 5, 10, 20, 50, 100, 200,
500, 1000\}$) produces a monotonically rightward-shifting family of sigmoid curves; the
inflection point $\beta^{*}(\rho)$ increases from 3.85 ($\rho=1$) to 10.22 ($\rho=1000$),
consistent with $\beta^{*} \propto \log K_{\mathrm{eff}}$ predicted by the
mean-field argument in Section~\ref{sec:theory}.
